\section{Conclusions and Future Work}
\label{sec:conclusions}

MouseClaimBench is an evidence-constrained \ac{kbs} for authorizing computational mouse-brain claims under an explicit policy. It externalizes a versioned profile, preserves source-scale facts and provenance, applies prioritized non-compensatory rules, and returns a machine-readable proof trace. Prediction, reproduction, topology, direction, local observational structure--function association, mechanism, causality, and complete entity-specific digital-twin scope remain distinct targets.

Exhaustive state tests establish deterministic rule behaviour for the mechanistic claim. In the known-truth benchmark, the policy obtains an \ac{fpr} of $0.019$, compared with $0.050$ for an equal-weight compensatory baseline and $0.271$ for a prediction shortcut. Its corresponding \ac{fnr} is $0.070$, compared with $0.041$ and $0.007$. Four real cases produce 40 bounded author-generated migration decisions that are reproduced with complete traces. External controls further show why retrieval cannot replace support inference and association cannot replace causal direction.

These results do not establish biological truth or external content validity. The profile is author-proposed, the four real cases are not a representative sample, and the \ac{microns} endpoint is internal to one volume. No case supports external replication, causal whole-brain identification, or a complete entity-specific digital twin. Future work must freeze the profile before prospective evaluation on independently selected studies, obtain expert assessment of its relations and proof traces, and test structure--function evidence across animals with network-aware inference. Another domain profile is justified only when its vocabulary and validation cases are independently available.

The conclusion is deliberately narrow. Accurate prediction and reproducible execution do not authorize mechanistic, causal, or digital-twin wording by themselves. MouseClaimBench represents this boundary as inspectable knowledge and records why each supported or restricted conclusion was reached.
